%!TEX root = ../dokumentation.tex

\chapter*{\abkuerzungstitel}
% Im Text: \ac{API} schreibt beim ersten Mal die Langform (als Fußnote), danach
% nur die Abkürzung. Gedruckt werden nur Abkürzungen, die im Text per \ac{}
% vorkommen (printonlyused); ohne jedes \ac{} bleibt die Liste leer.
% [WYSIWYG] = längste Abkürzung, bestimmt die Spaltenbreite.
\begin{acronym}[WYSIWYG]
\setlength{\itemsep}{-\parsep}

\acro{API}{Application Programming Interface}
\acro{DHBW}{Duale Hochschule Baden-Württemberg}
\acro{PDF}{Portable Document Format}
\end{acronym}
